% 原书第 33 页无编号示意图；保留算法库、单折训练/验证划分与 V 折面板。
\begin{tikzpicture}[x=1.2cm,y=1cm,line width=.45pt,font=\small]
 \draw (0,.80) rectangle (1.50,3.10);
 \node[align=center,font=\small] at (.75,1.95) {glm\\glmnet\\$\vdots$\\nnet};
 \node at (.75,.55) {算法};

 % 左侧放大的第 1 折；灰色为验证集。
 \fill[gray!35] (2.30,.25) rectangle (3.15,.95);
 \draw (2.30,.25) rectangle (3.15,3.80);
 \draw (2.30,.95)--(3.15,.95);
 \draw (2.30,3.10)--(3.15,3.10);
 \node at (2.725,3.45) {$1$};
 \node at (2.725,1.95) {$\vdots$};
 \node at (2.725,.60) {$V$};
 \node at (2.725,-.12) {第 1 折};
 \draw[<->,>=Stealth] (3.48,1.00)--(3.48,3.75);
 \draw[<->,>=Stealth] (3.48,.27)--(3.48,.92);
 \draw (3.33,3.80)--(3.63,3.80);
 \draw (3.33,.95)--(3.63,.95);
 \draw (3.33,.25)--(3.63,.25);
 \node[align=center,anchor=west] at (3.72,2.4) {训练集};
 \node[align=center,anchor=west] at (3.72,.60) {验证集};

 % 右侧 V 折交叉验证面板；不虚构省略的中间折。
 \fill[gray!35] (5.40,.35) rectangle (6.15,1.05);
 \draw (5.40,.35) rectangle (6.15,3.70);
 \draw (5.40,1.05)--(6.15,1.05);
 \draw (5.40,3.00)--(6.15,3.00);
 \node at (5.775,3.35) {$1$};
 \node at (5.775,2.03) {$\vdots$};
 \node at (5.775,.70) {$V$};
 \node at (5.775,-.02) {第 1 折};

 \fill[gray!35] (6.40,1.05) rectangle (7.15,1.75);
 \draw (6.40,.35) rectangle (7.15,3.70);
 \draw (6.40,1.05)--(7.15,1.05);
 \draw (6.40,1.75)--(7.15,1.75);
 \draw (6.40,3.00)--(7.15,3.00);
 \node at (6.775,3.35) {$1$};
 \node at (6.775,2.35) {$\vdots$};
 \node at (6.775,.70) {$V$};
 \node at (6.775,-.02) {第 2 折};

 \fill[gray!35] (7.40,1.75) rectangle (8.15,2.45);
 \draw (7.40,.35) rectangle (8.15,3.70);
 \draw (7.40,1.05)--(8.15,1.05);
 \draw (7.40,1.75)--(8.15,1.75);
 \draw (7.40,2.45)--(8.15,2.45);
 \draw (7.40,3.00)--(8.15,3.00);
 \node at (7.775,3.35) {$1$};
 \node at (7.775,2.70) {$\vdots$};
 \node at (7.775,1.40) {$\vdots$};
 \node at (7.775,.70) {$V$};
 \node at (7.775,-.02) {第 3 折};

 \draw (8.40,.35) rectangle (9.15,3.70);
 \node at (8.775,2.00) {$\vdots$};
 \node at (8.775,-.02) {$\cdots$};

 \fill[gray!35] (9.40,3.00) rectangle (10.15,3.70);
 \draw (9.40,.35) rectangle (10.15,3.70);
 \draw (9.40,1.05)--(10.15,1.05);
 \draw (9.40,3.00)--(10.15,3.00);
 \node at (9.775,3.35) {$1$};
 \node at (9.775,2.03) {$\vdots$};
 \node at (9.775,.70) {$V$};
 \node at (9.775,-.02) {第 $V$ 折};
\end{tikzpicture}
